\section*{1. Tóm tắt}

Bài luận này nghiên cứu cải tiến hệ thống truyền thông hỗ trợ học viên tại Trung tâm Công nghệ Giáo dục (daotao.ai), Đại học Bách khoa Hà Nội, trong bối cảnh đào tạo số quy mô lớn và hỗ trợ Đề án 06 của Chính phủ.

Dựa trên bốn lý thuyết truyền thông cốt lõi gồm mô hình Shannon–Weaver, mô hình giao dịch, thuyết phong phú truyền thông và thuyết tải nhận thức, nghiên cứu chẩn đoán các điểm đứt gãy tương tác và đề xuất hai giải pháp cải tiến thực tế: tích hợp chatbot tự động trợ giúp trên website để giải quyết các sự cố thường gặp như lỗi đăng nhập hay nút đăng ký và thiết kế bộ mẫu thư điện tử trực quan tinh giản thông tin. Đồng thời, trung tâm áp dụng quy trình bàn giao công việc qua bảng quản trị nội bộ nhằm hạn chế lỗi kỹ thuật học liệu.

Kết quả thử nghiệm thực tế chứng minh hiệu quả của các giải pháp đề xuất: nhóm áp dụng cải tiến mới ghi nhận sự sụt giảm đáng kể số cuộc gọi cần hỗ trợ kỹ thuật trực tiếp, nâng cao tỷ lệ học viên tự thao tác đăng nhập thành công ngay trong ngày đầu tiên, đồng thời cải thiện tỷ lệ học viên hoàn thành bài kiểm tra điều kiện của khóa học. Nghiên cứu khẳng định giá trị thực tế của việc ứng dụng khoa học truyền thông vào tối ưu hóa vận hành đào tạo trực tuyến.

\vspace{0.5cm}
\noindent \textbf{Từ khóa:} Khoa học truyền thông, Công nghệ giáo dục, Trí tuệ nhân tạo, Tải nhận thức, daotao.ai, Đề án 06.
